% !TeX program = xelatex
% UTF-8 transcription of all three images in Google Drive Notebook/hemolytic anemia.
% pic01.jpg = 20261016_225127862.jpg; source page 1.
% pic02.jpg = 20261016_225134441.jpg; source page 2.
% pic03.jpg = 20261016_225140709.jpg; source page 3.
% Source wording and laboratory statements are retained without medical correction.
% Standalone XeLaTeX document, also suitable for input into a XePersian document.
\newif\ifHemolyticStandalone
\ifdefined\settextfont
  \HemolyticStandalonefalse
\else
  \HemolyticStandalonetrue
\fi
\ifHemolyticStandalone
  \documentclass[11pt,a4paper]{article}
  \usepackage[margin=18mm]{geometry}
  \usepackage{amsmath,amssymb,array}
  \usepackage{xepersian}
  \IfFontExistsTF{Amiri}{\settextfont{Amiri}}{\settextfont{FreeSerif}}
  \setlatintextfont{FreeSerif}
  \begin{document}
\fi
\begingroup
\setlength{\parindent}{0pt}
\setlength{\parskip}{5pt}
\emergencystretch=3em
\newcommand{\HemolyticEn}[1]{\lr{#1}}
\newcommand{\HemolyticArrow}{\ensuremath{\longleftarrow}}

\begin{latin}
\textbf{Harrison}\hfill Subject: Hemolytic Anemias -- 100\hfill Date: /\quad/
\end{latin}
\section*{دسته‌بندی آنمی‌های همولیتیک}
\begin{center}
\renewcommand{\arraystretch}{1.5}
\begin{tabular}{|p{0.12\linewidth}|p{0.36\linewidth}|p{0.38\linewidth}|}
\hline
 & \textbf{فاکتورهای خارج سلولی} & \textbf{نقص داخل سلولی} \\
\hline
ارثی & \HemolyticEn{Familial (atypical) hemolytic uremic syndrome}
& هموگلوبینوپاتی‌ها؛ آنزیموپاتی‌ها؛ نقص غشا ـ سایتواسکلتون \\
\hline
اکتسابی & خودایمنی؛ عفونت‌ها؛ داروها / سموم؛ میکروآنژیوپاتی‌ها
& \HemolyticEn{PNH} \\
\hline
\end{tabular}
\end{center}

\subsection*{آنمی‌های همولیتیک بر اساس عمده‌ترین محل همولیز}
\begin{itemize}
\item \HemolyticEn{Intravascular}
\item \HemolyticEn{Extravascular}
\end{itemize}

\subsection*{آنمی همولیتیک}
\begin{itemize}
\item \HemolyticEn{Abrupt} \HemolyticArrow{} \HemolyticEn{G6PDD, AIHA}.
\item \HemolyticEn{Gradual} \HemolyticArrow{} \HemolyticEn{CAD, HS}.
\end{itemize}
\begin{latin}
CAD $\longrightarrow$ Cold Aglutinin Disease
\end{latin}

\subsection*{افتراق آنمی همولیتیک از سایر آنمی‌ها}
علائم و نشانه‌های همولیز.
\begin{itemize}
\item علامت اصلی در بالین \HemolyticArrow{} زردی.
\item سایر علائم \HemolyticArrow{} تغییر رنگ ادرار.
\end{itemize}
\vfill
\begin{center}\HemolyticEn{YASHA}\qquad ۱\end{center}

\clearpage
\begin{latin}
Subject: 100\hfill Date: /\quad/
\end{latin}
\section*{آنمی همولیتیک ـ صفحهٔ ۲}
\begin{itemize}
\item در بسیاری از موارد \HemolyticEn{HA} \HemolyticArrow{} اسپلنومگالی.
\item در برخی موارد \HemolyticEn{HA} \HemolyticArrow{} هپاتومگالی.
\item در همهٔ موارد \HemolyticEn{HA} شدید مادرزادی \HemolyticArrow{}
تغییرات اسکلتی ـ به علت فعالیت بیش از حد \HemolyticEn{BM}.
\end{itemize}

\begin{latin}
\subsection*{Patient with HA}
\begin{enumerate}
\renewcommand{\labelenumi}{\alph{enumi})}
\item G.A $\longrightarrow$ Jaundice, Pallor
\item Ph/E $\longrightarrow$ Splenomegally, Bossing of skull (if severe)
\item Hb $\longrightarrow$ N/L to very low
\item MCV, MCH $\longrightarrow$ Usually $\uparrow$
\item Retic. $\longrightarrow$ Usually $\uparrow$
\item Bilirubin $\longrightarrow$ Almost always $\uparrow$ (mostly unconjugated)
\item LDH $\longrightarrow$ $\uparrow$ (up to 10X normal in intravascular hemolysis)
\item Haptoglobin $\longrightarrow$ Reduced to absent if hemolysis is at least in part intravascular
\end{enumerate}
\end{latin}

\subsection*{همولیزی که عمدتاً خارج عروقی باشد}
\begin{itemize}
\item افزایش بیلی‌روبین غیرکنژوگهٔ سرم.
\item افزایش \HemolyticEn{AST} سرم.
\item افزایش اوروبیلینوژن ادرار و مدفوع.
\end{itemize}
\vfill
\begin{center}\HemolyticEn{YASHA}\qquad ۲\end{center}

\clearpage
\begin{latin}
Subject: 100\hfill Date: /\quad/
\end{latin}
\section*{همولیزی که عمدتاً داخل عروقی باشد}
\begin{itemize}
\item هموگلوبین‌اوری (اغلب همراه با هموسیدرین‌اوری).
\item هموگلوبین آزاد در سرم.
\item افزایش \HemolyticEn{LDH} سرم.
\item کاهش هاپتوگلوبین سرم.
\item بیلی‌روبین سرم نرمال یا اندکی افزایش‌یافته.
\end{itemize}

افزایش تعداد رتیکولوسیت‌ها در \HemolyticEn{HA} \HemolyticArrow{}
افزایش \HemolyticEn{MCV}.

\HemolyticEn{BMA} در تشخیص \HemolyticEn{HA} جایگاهی ندارد.
\vfill
\begin{center}\HemolyticEn{YASHA}\qquad ۳\end{center}
\endgroup
\ifHemolyticStandalone
  \expandafter\enddocument
\fi
